\documentclass[11pt,a4paper]{article}

\usepackage[utf8]{inputenc}
\usepackage[T1]{fontenc}
\usepackage[ngerman]{babel}
\usepackage{lmodern}
\usepackage{amsmath,amssymb}
\usepackage{graphicx}
\usepackage{booktabs,tabularx,array}
\usepackage[table]{xcolor}
\usepackage{tikz}
\usetikzlibrary{arrows.meta,positioning}
\usepackage[most]{tcolorbox}
\usepackage{enumitem}
\usepackage{microtype}
\usepackage{hyperref}
\usepackage{listings}
\usepackage[a4paper,left=18mm,right=18mm,top=15mm,bottom=15mm,headheight=0pt]{geometry}

\definecolor{navy}{HTML}{174A63}
\definecolor{blue}{HTML}{2E7698}
\definecolor{lightblue}{HTML}{EAF3F7}
\definecolor{lightgray}{HTML}{F2F3F4}
\definecolor{midgray}{HTML}{66727A}
\definecolor{green}{HTML}{477A63}
\definecolor{orange}{HTML}{B96A2B}
\definecolor{safe}{HTML}{287B61}
\definecolor{safebg}{HTML}{E7F6EF}
\definecolor{danger}{HTML}{AF3F47}
\definecolor{dangerbg}{HTML}{FFF1F1}

\hypersetup{colorlinks=true,linkcolor=navy,urlcolor=blue,pdftitle={Sicherungsblatt Aufgabe 6 -- Perzeptron implementieren, einfache Variante}}
\setlength{\parindent}{0pt}
\setlength{\parskip}{4pt}
\setlist[itemize]{leftmargin=5mm,itemsep=1pt,topsep=1pt}
\renewcommand{\familydefault}{\sfdefault}
\pagestyle{empty}

\newtcolorbox{definition}[1][]{enhanced,colback=lightblue,colframe=blue,
  boxrule=0.8pt,arc=1.5mm,left=3mm,right=3mm,top=1.5mm,bottom=1.5mm,
  fonttitle=\bfseries,title=Definition,#1}
\newtcolorbox{merke}[1][]{enhanced,colback=lightgray,colframe=navy,
  boxrule=0.9pt,arc=1.5mm,left=3mm,right=3mm,top=1.5mm,bottom=1.5mm,
  fonttitle=\bfseries,title=Merke,#1}
\newtcolorbox{beispiel}[1][]{enhanced,colback=white,colframe=midgray,
  boxrule=0.7pt,arc=1.5mm,left=3mm,right=3mm,top=1.5mm,bottom=1.5mm,
  fonttitle=\bfseries,title=Beispiel,#1}

\newcommand{\sectionline}[1]{%
  \vspace{0.5mm}{\Large\bfseries\color{navy}#1}\par\vspace{0.4mm}%
  {\color{blue}\rule{\linewidth}{0.8pt}}\vspace{0.7mm}}
\newcommand{\sheettitle}[2]{%
  \begin{tcolorbox}[colback=navy,colframe=navy,arc=2mm,boxrule=0pt,left=5mm,right=5mm,top=2.5mm,bottom=2.5mm]
    {\color{white}\fontsize{19}{22}\selectfont\bfseries #1}\par
    \vspace{0.5mm}{\color{white}\large #2}
  \end{tcolorbox}}
\lstdefinestyle{java-sheet}{
  basicstyle=\ttfamily\footnotesize,
  columns=fullflexible,keepspaces=true,
  showstringspaces=false,aboveskip=0pt,belowskip=0pt
}

\begin{document}

\sheettitle{Perzeptron implementieren}{Einfache Variante · Sicherungsblatt zu Aufgabe 6 · Informatik 11}

\sectionline{1. Klassifizieren: berechnen und entscheiden}
\begin{center}
\begin{tikzpicture}[>=Latex,font=\small,
  flowbox/.style={draw=blue,fill=lightblue,rounded corners=1.5mm,minimum height=10mm,align=center,inner sep=2mm},
  result/.style={draw=safe,fill=safebg,rounded corners=1.5mm,minimum height=10mm,align=center,inner sep=2mm}]
  \node[flowbox,minimum width=27mm] (input) at (0,0) {$x_1,\;x_2$\\Eingaben};
  \node[flowbox,minimum width=37mm] (sum) at (4.1,0) {$a=w_1x_1+w_2x_2$\\gewichtete Summe};
  \node[flowbox,minimum width=28mm] (test) at (8.6,0) {$a\ge\theta$\,?\\Vergleich};
  \node[result,minimum width=29mm] (out) at (12.4,0) {ja: \textbf{1}\\nein: \textbf{0}};
  \draw[->,thick,navy] (input) -- (sum);
  \draw[->,thick,navy] (sum) -- (test);
  \draw[->,thick,navy] (test) -- (out);
\end{tikzpicture}\par
{\small\color{midgray}\texttt{punktKlassifizieren} gibt nur 0 oder 1 zurück. Gewichte und \texttt{theta} bleiben dabei gleich.}
\end{center}

\sectionline{2. Die drei Ergänzungen in \texttt{trainieren}}
\begin{minipage}[t]{0.49\linewidth}
\vspace{0pt}
\begin{definition}[title=Schritt 2 · Summe als \texttt{double}]
\begin{lstlisting}[style=java-sheet]
double gewichtete_Summe =
    w_1 * punkt.gibEingabewertX1()
  + w_2 * punkt.gibEingabewertX2();
\end{lstlisting}
\end{definition}
\end{minipage}\hfill
\begin{minipage}[t]{0.49\linewidth}
\vspace{0pt}
\begin{definition}[title=Schritt 3 · Ausgabe als \texttt{int}]
\begin{lstlisting}[style=java-sheet]
if (gewichtete_Summe >= theta) {
    berechnete_ausgabe = 1;
} else {
    berechnete_ausgabe = 0;
}
\end{lstlisting}
\end{definition}
\end{minipage}

\vspace{0.5mm}
\begin{center}
\colorbox{lightgray}{\parbox{0.94\linewidth}{\centering\small
\texttt{int delta = punkt.gibLabel() - berechnete\_ausgabe;}
\quad $\Longrightarrow$ \quad $-1,\;0$ oder $+1$}}
\end{center}

\begin{minipage}[t]{0.49\linewidth}
\begin{beispiel}[title=Schritt 4 · \textcolor{safe}{$\delta=+1$}]
\begin{lstlisting}[style=java-sheet]
if (delta == 1) {
    w_1 = w_1 + lernrate *
        punkt.gibEingabewertX1();
    w_2 = w_2 + lernrate *
        punkt.gibEingabewertX2();
    theta = theta - lernrate;
}
\end{lstlisting}
\end{beispiel}
\end{minipage}\hfill
\begin{minipage}[t]{0.49\linewidth}
\begin{beispiel}[title=Schritt 4 · \textcolor{danger}{$\delta=-1$}]
\begin{lstlisting}[style=java-sheet]
else if (delta == -1) {
    w_1 = w_1 - lernrate *
        punkt.gibEingabewertX1();
    w_2 = w_2 - lernrate *
        punkt.gibEingabewertX2();
    theta = theta + lernrate;
}
\end{lstlisting}
\end{beispiel}
\end{minipage}
\vspace{0.3mm}\par
{\small\textbf{Bei $\delta=0$:} Keine der beiden Bedingungen trifft zu; \texttt{w\_1}, \texttt{w\_2} und \texttt{theta} bleiben unverändert.}

\sectionline{3. Programmlauf mit den vier Trainingspunkten}
{\small Start: $\boxed{(w_1,w_2,\theta)=(1,1,1)}$, Lernrate $=1$. Die Punkte werden in dieser Reihenfolge trainiert:}
\vspace{0.6mm}\par
\begin{center}
\renewcommand{\arraystretch}{1.2}
\begin{tabularx}{0.94\linewidth}{@{}>{\centering\arraybackslash}p{18mm}>{\centering\arraybackslash}p{21mm}>{\centering\arraybackslash}p{15mm}X>{\centering\arraybackslash}p{33mm}@{}}
\toprule
Punkt $(x_1|x_2)$ & Label $t$ & $\delta$ & Wirkung & $(w_1,w_2,\theta)$ danach\\
\midrule
$(2|2)$ & 1 & 0 & keine Änderung & $(1,1,1)$\\
\rowcolor{dangerbg}$(4|1)$ & 0 & $-1$ & $w_1,w_2$ sinken; $\theta$ steigt & $(-3,0,2)$\\
$(6|2)$ & 0 & 0 & keine Änderung & $(-3,0,2)$\\
\rowcolor{safebg}$(0|4)$ & 1 & $+1$ & $w_2$ steigt; $\theta$ sinkt & $(-3,4,1)$\\
\bottomrule
\end{tabularx}
\end{center}
\begin{merke}[title=Vorher $\longrightarrow$ nachher]
{\small $(4|1)$: zu Beginn $1\cdot4+1\cdot1=5\ge1\Rightarrow\boxed{1}$;\quad
nach dem Training $-3\cdot4+4\cdot1=-8<1\Rightarrow\boxed{0}$.\\
\texttt{double}: Eingaben, Gewichte, Lernrate, Summe · \texttt{int}: Label, Ausgabe, \texttt{delta}.}
\end{merke}

\end{document}
